% Book Two Pass 3: source-derived original scientific plate.
\clearpage
\section*{A 3D view of a 2D numerical result}
\noindent A planar scalar field can be lifted into three dimensions to reveal structure. Here the vertical axis represents estimated geodesic distance, not physical terrain elevation. The second panel shows signed error directly over the actual triangle mesh.

\begin{figure}[!htbp]
\centering
\includegraphics[width=\linewidth]{book02-figures/rendered/figures/geodesic-surface-3d.pdf}
\caption{Three-dimensional height field of the reconstructed distance and corresponding signed-error map with original mesh edges. Both are derived from the same computed numerical solution and exact flat-domain reference.}
\label{fig:book02-atlas-geodesic-surface-3d}
\end{figure}

\noindent\textbf{Reconstruction experiment.} Draw the analytical distance as a second surface, then plot their difference. Change the camera angle and explain which numerical defects become more or less visually apparent.

\noindent\emph{Reproducibility:} The original figure, SVG, print PDF, PNG,
exact numerical CSV data and machine-readable provenance report are produced
by \texttt{tools/labs/book02\_visual\_atlas.py}. These are independently
written educational experiments, not imported research-paper figures.
